\documentclass[11pt,a4paper]{article}
\usepackage[T1]{fontenc}
\usepackage[utf8]{inputenc}
\usepackage[ngerman]{babel}
\usepackage[margin=2.5cm]{geometry}
\usepackage[hidelinks]{hyperref}
\setlength{\parindent}{0pt}
\setlength{\emergencystretch}{2em}
\begin{document}
\begin{center}
  {\Large\textbf{Webtechnologien -- Übungsblatt 1}}
\end{center}
\section*{Aufgabe 4: Bildergalerie}
\textbf{Aufgabe:} Eine Startseite und drei Bildseiten mit den
vorgegebenen Bildern und Texten erstellen. Die Bildseiten sollen nach
sieben Sekunden automatisch zum nächsten Bild wechseln.

\medskip
\textbf{Ergebnis:}
\begin{itemize}
  \item \texttt{index.html}: Überschrift \glqq Image Gallery\grqq{} und drei Vorschaubilder.
    Der Bereich hat die ID \texttt{gallery-section}; die Bilder verlinken die Detailseiten.
  \item \texttt{img1.html}, \texttt{img2.html}, \texttt{img3.html}:
    Bildtitel, Bild und vorgegebene Beschreibung.
  \item Bildtitel beim Überfahren der Vorschaubilder;
    \glqq Open image [Titel] in full\grqq{} auf den Bildseiten.
  \item Klick auf das große Bild öffnet die JPEG-Datei in einem neuen Tab.
  \item Links zu vorherigem Bild, nächstem Bild und Startseite.
  \item Automatischer Wechsel per Meta-Refresh nach sieben Sekunden:
    Bild 1 $\rightarrow$ Bild 2 $\rightarrow$ Bild 3 $\rightarrow$ Bild 1.
\end{itemize}
\textbf{alt-Attribute:} Beschreiben den Bildinhalt für Screenreader
und bei nicht geladenen Bildern.

\medskip
\textbf{Test:} Alle vier HTML-Seiten und sechs Bilddateien sind über
Apache mit Status \texttt{200 OK} erreichbar.
Startseite: \url{http://localhost/webtech-aufgabe4/}.
\medskip
{\small \textbf{Hilfsmittel:} Codex für Erstellung, Prüfung und Dokumentation.}
\end{document}
